\chapter{糖酵解：把葡萄糖一点点拆开}

\begin{kaichang}
我们的老朋友面包又来了。第 46 章里，我们拆开过它的淀粉——那是一长串葡萄糖手拉手连成的大分子。这一次，请你关注面包的另一样东西：那些大大小小的孔。

把一团刚和好的面放在温暖的地方，一两个小时后它会胀到原来两倍大，按一下，软绵绵地弹回来；撕开看，里面全是蜂窝。烤成面包以后，孔就定型了。这些孔是谁吹出来的？答案你知道：酵母。可酵母只是一个直径不到头发丝十分之一的单细胞真菌，它没有嘴，没有牙，没有肺，更没有打气筒。它只是“泡”在面团里，就把糖变成了气体——顺便还会产生一点酒精（别担心，烘烤时几乎全挥发了）。

请在纸上写下你的猜测：酵母细胞内部到底发生了什么，才能一边“吃糖”一边吹出气泡？是糖在细胞里“烧起来”了吗？如果是烧，为什么面团没有烫手、没有火苗？这一章结束时，我们回来对答案。
\end{kaichang}

\section{面团里看不见的热闹}

先把能观察到的东西摆出来。

温水化开一小撮酵母，加一勺糖，过不了多久，液面浮起泡沫，能闻到淡淡的酒香和酸味；摸一摸容器壁，微微发温。这说明至少有三件事在发生：\textbf{有气体生成}（泡沫）、\textbf{有新物质生成}（酒香、酸味）、\textbf{有能量以热的形式散出}（发温）。气体用澄清石灰水一检验，变浑浊——是二氧化碳，和你呼出的是同一种分子（第 32 章）。

再看几个似曾相识的场面。牛奶放久了变成酸奶，酸味的来源是一种叫乳酸的物质；剧烈运动之后你大口喘气，肌肉又酸又沉；酸菜坛子、葡萄酒缸、馒头屉——这些风马牛不相及的东西，背后站着同一群微生物，做着同一桩生意：\textbf{在没有氧气（或氧气不够）的地方，从糖身上榨取能量}。

现在把镜头推近，一直推到酵母细胞的内部。氧气进不来或者不够用，可细胞仍然要支付各种“账单”：合成原料、维持膜两侧的浓度差（第 51 章）、修补零件。钱从哪里来？细胞没有火炉，不能点火——蛋白质零件在火焰温度下早就散架了（第 48 章讲过的变性）。它唯一的本钱，是第 46 章认识的葡萄糖分子：六个碳原子组成的环，身上挂满了羟基。

接下来发生的，就是本章的主角：\textbf{糖酵解}——“糖”指葡萄糖，“酵解”就是“把它拆解开来”。顾名思义，这是一条把一个六碳糖拆成两个三碳碎片的反应流水线。

\section{十步接力：先存钱，才能取钱}

糖酵解一共十步反应，由十种酶接力完成（酶的本事见第 49 章），全部发生在细胞质里，不需要氧气，也不需要任何细胞器——连没有线粒体的红细胞、连比酵母古老得多的细菌，都会这一套。它是地球上分布最广的代谢通路，从细菌到你，用的是几乎一模一样的十步。

十步不用背，但整条流水线分两个阶段的“生意经”必须看懂：

\begin{itemize}[itemsep=2pt]
\item \textbf{投资阶段（前五步）}：细胞先\textbf{倒贴}两分子 ATP。ATP 把磷酸基团“贴”到葡萄糖身上（第 50 章讲过，磷酸基团转移能改变分子的反应性），葡萄糖被逐步改造成一个容易从中间掰开的形状，然后“咔嚓”裂成两个三碳的小分子。这一步很像做生意先垫本钱。
\item \textbf{收益阶段（后五步）}：两个三碳分子各自再经历几步改造，一边被拆掉零件，一边“找零”——每一步放出的小额能量被当场用来合成 ATP，同时从碳骨架上剥下电子，交给电子载体 \chem{NAD^{+}}，变成 NADH。最终每个三碳分子变成一个三碳的丙酮酸。
\end{itemize}

算总账：投入 2 个 ATP，产出 4 个 ATP，\textbf{净赚 2 个 ATP}；另外进账 2 个 NADH（里面装着从葡萄糖身上剥下来的高能电子）；剩下两个丙酮酸分子，装着还没拆完的大部分碳骨架。

\begin{center}
\begin{tikzpicture}[scale=1.0]
\node[draw,rectangle,minimum width=2.2cm,minimum height=0.8cm] (a) at (0,0) {葡萄糖};
\node[draw,rectangle,minimum width=3.0cm,minimum height=0.8cm] (b) at (4.4,0) {六碳中间物};
\node[draw,rectangle,minimum width=3.0cm,minimum height=0.8cm] (c) at (8.8,0) {两个三碳糖};
\node[draw,rectangle,minimum width=2.4cm,minimum height=0.8cm] (d) at (12.6,0) {2 丙酮酸};
\draw[->,thick] (a) -- (b) node[midway,above]{\small 投资 $2\,\mathrm{ATP}$};
\draw[->,thick] (b) -- (c) node[midway,above]{\small 从中间裂开};
\draw[->,thick] (c) -- (d) node[midway,above]{\small 收益 $4\,\mathrm{ATP}$};
\node at (10.4,-1.0) {\small 同时生成 $2\,\mathrm{NADH}$（装着高能电子）};
\end{tikzpicture}
\end{center}

这幅图是人为的表示法：方框和箭头只标出物质和“钱”的进出，不代表分子的真实形状，实际的十步比这细得多。

\section{流水线上的四个关键工位}

十步反应的名字不用记，但其中有四个“工位”值得停下来看一眼，因为它们分别回答了四个真实的问题。

\textbf{第一个问题：葡萄糖凭什么留在细胞里？}葡萄糖顺着浓度梯度进入细胞后，照理说也能顺原路溜出去。第一步酶解决的办法是“盖章”：从 ATP 身上掰下一个磷酸基团，贴在葡萄糖的第 6 号碳上。贴了磷酸的葡萄糖带上了负电荷，膜上那些只认中性分子的通道再也认不出它——货物一进仓库就换了包装，退不了货了。这就是投资阶段花掉的第一个 ATP 买来的东西：\textbf{原料的独占权}。

\textbf{第二个问题：六碳骨架从哪儿下刀？}直接投资阶段最贵的一刀：第三个工位再添一个磷酸基团，把分子改造成一根“两边都被拽住”的双磷酸链条，下一个酶正好从正中间把它劈成两个三碳糖。先绷紧，再剪断——垫钱不是浪费，是把分子摆到“一掰就开”的姿势上。

\textbf{第三个问题：电子从哪里剥下来？}收益阶段的第一步是全程唯一的氧化反应（第 28 章：氧化的本质是失电子）。酶把三碳糖骨架上的一个醛基氧化，剥下的电子连同一个质子一起，被等在旁边的 \chem{NAD^{+}} 接走，变成 NADH。葡萄糖储存的能量，相当一部分就藏在这对被剥走的电子里——\textbf{NADH 是本章真正的“高息存单”}，可惜它要到第 53 章才能兑现。

\textbf{第四个问题：ATP 在哪里造出来？}收益阶段有两步，中间产物身上的磷酸基团被直接“过户”给 ADP，当面合成 ATP。注意这里没有氧气、没有线粒体什么事：磷酸基团从一个底物分子手里转交给 ADP，这种方式叫\textbf{底物水平磷酸化}。它产 ATP 的速度很快，但产量有限——每分子葡萄糖只有这区区 4 次“当面转账”。人体绝大部分 ATP 其实来自另一种更壮观的方式，那是下一章的主角。

\section{谁在给流水线踩刹车}

现在回答大纲里最容易被忽略、却最接近生命本质的那个问题：这条流水线怎样受控制？

细胞不可能一直全速拆糖——糖是宝贵的储备，ATP 用不完时还在猛烧，等于开着水龙头灌一个已经满了的池子。控制的办法不是“细胞做个决定”，而是写进分子里的化学反馈：糖酵解第三个工位那台酶（就是添第二个磷酸基团的那台），恰好是全流水线最慢、最不可逆的一道闸门，而它身上的别构位点（第 49 章讲过，活性位点之外的“遥控开关”）能直接感知细胞的能量状态——ATP 浓度高，意味着“钱够花了”，ATP 就结合上去把这台酶\textbf{关掉}；相反，ATP 被花掉后的产物 AMP 积累起来，说明“账上没钱了”，AMP 结合上去把酶\textbf{重新打开}。闸门酶同时听着“收入”和“欠条”两边的声音，整条流水线的速度就跟着能量收支自动浮沉。

这正是五个贯穿问题里第四问“变化怎样受控制”的一次实战：控制不需要指挥官，只需要酶分子对不同小分子的亲和力差异。至于进食、饥饿、运动时，激素怎样从全身尺度上再给这套局部反馈加一层指挥——胰岛素和胰高血糖素的故事，第 55 章会展开。

\section{为什么不一把火烧掉？}

你可能会嘀咕：费这么大劲，十步，才净赚 2 个 ATP？把葡萄糖直接和氧气烧掉，放出的能量不是多得多吗？

确实多得多——但“多”恰恰是问题。回忆第 27 章的道理：能量释放是一回事，\textbf{能量以什么速率、什么形式释放}是另一回事。葡萄糖和氧气的反应在纸面上放能巨大，可一旦发生就是燃烧——能量在百万分之一秒内以热的形式一涌而出。对细胞来说，这就像领工资时有人把一整年的钞票一次性点着塞给你：数额可观，一分也花不出去。

细胞的真正需求很具体：每一次“消费”——收缩一下肌肉、泵一批离子、合成一段蛋白质——需要的都只是几十个 ATP 量级的小额支付（第 50 章）。所以聪明的策略不是“放能最多”，而是\textbf{把一大笔能量拆成许多小份，每一份在释放的瞬间被另一个分子当场接住}。糖酵解的十步就是十个小台阶：每步只放一点点自由能，每步都有专门的酶守在台阶口，把这点能量转手塞进 ATP 或 NADH 里。火是把钱烧了取暖，酶是当面点钞。

\begin{qiaoliang}
这笔账可以定量算一下。实验测得：1 摩尔葡萄糖（\ud{180}{g}）在氧气中完全燃烧成二氧化碳和水，自由能变化约 \ud{-2870}{kJ/mol}。而在细胞条件下，水解 1 摩尔 ATP 大约能释放 \ud{30}{kJ} 可用的能量。

糖酵解每摩尔葡萄糖净存 2 摩尔 ATP，存下约 \ud{60}{kJ}——只占 2870 的百分之二左右！其余的能量去哪了？没有消失（能量守恒可不是开玩笑）：大部分还锁在两个丙酮酸和两分子 NADH 的化学键里，一小部分以热的形式散掉了（这就是面团发温的原因之一）。

换句话说，糖酵解根本不是一次完整的“取款”，只是从葡萄糖这张大钞票上\textbf{先剪下来两个零头}应急。整张钞票怎么兑开？那是第 53 章的故事。顺便再体会一下“小额支付”的规模：你全身的细胞每天合成又消耗掉的 ATP，总质量大约相当于你的体重——每一分子 ATP 平均每分钟就要被“花掉再挣回”好几轮。没有这套接力拆糖的本事，你连读完这句话的能量都付不起。
\end{qiaoliang}

\section{符号登场：把十步压缩成一行}

现在我们才请出符号。把十步的所有中间产物略去，糖酵解的总账目可以写成：

\[\chem{C_6H_{12}O_6} + \chem{2ADP} + \chem{2P_{i}} + \chem{2NAD^{+}} \rightarrow \chem{2C_3H_4O_3} + \chem{2ATP} + \chem{2NADH} + \chem{2H^{+}} + \chem{2H_2O}\]

从左读到右：一分子葡萄糖（\chem{C_6H_{12}O_6}），消耗两分子 ADP、两分子磷酸（\chem{P_{i}}，游离的磷酸根）和两分子 \chem{NAD^{+}}；生成两分子丙酮酸（\chem{C_3H_4O_3}）、两分子 ATP 和两分子 NADH。请特别盯一眼 \chem{NAD^{+}} 这个角色：它出现在反应物里。这意味着——\textbf{细胞里的 \chem{NAD^{+}} 一旦用光，整条流水线立刻停工}。记住这句话，它是理解下半章的钥匙。

\section{\chem{NAD^{+}} 危机与发酵的救场}

现在轮到那个被忽略的问题登场了。糖酵解每跑一轮，就要消耗一分子 \chem{NAD^{+}}（其实是两分子）来接收电子。而细胞里 \chem{NAD^{+}} 的存货很少，全部转化成大写字母也就能写几行——打个比方，它像收银台抽屉里的零钱，几分钟就会被找光。没有 \chem{NAD^{+}}，糖酵解这台机器就卡死在半路上，ATP 断供，细胞几分钟内就要“破产”。

在有氧气、有线粒体的细胞里，NADH 可以把电子“快递”给线粒体，卸完货变回 \chem{NAD^{+}}，循环再用（第 53 章细讲这条快递线）。可是面团深处的酵母、酸奶缸里的乳酸菌、你剧烈收缩时氧气送不及的肌肉细胞，都面临同一个处境：\textbf{快递线断了或运力不够，NADH 堆积，\chem{NAD^{+}} 告急}。

怎么办？这些细胞用了一招朴素的会计手段：让 NADH 把电子就地“倒”给糖酵解自己的产物。这一招叫\textbf{发酵}。最常见的两套做法：

\begin{itemize}[itemsep=2pt]
\item \textbf{乳酸发酵}：NADH 把电子直接塞给丙酮酸，丙酮酸接住电子就变成乳酸。你的肌肉细胞在缺氧冲刺时这么干；把牛奶变成酸奶、把白菜变成酸菜的乳酸菌也这么干——酸味就是乳酸的味道。
\item \textbf{酒精发酵}：酵母先把丙酮酸削掉一个碳（以一分子二氧化碳的形式放走——面包孔的来历！），剩下的两碳碎片接住 NADH 的电子，变成酒精。
\end{itemize}

看明白了吗？乳酸和酒精\textbf{不是发酵的目的，是发酵的副产品}。细胞真正要的东西只有一样：把 NADH 变回 \chem{NAD^{+}}，好让糖酵解继续转、继续产那可怜但救命的 2 个 ATP。所以别再把发酵看成“无氧呼吸的低级版本”——它不是一条产能通路，而是产能通路的\textbf{续命装置}。乳酸和酒精只是细胞倒掉的“电子垃圾”，只不过这堆垃圾碰巧酸得开胃、香得醉人，被人类做成了一门几千年的生意。

这里还有一个常常让人困惑的问题：运动后肌肉酸痛，是不是乳酸在作怪？

\begin{wuqu}
“运动后第二天肌肉酸痛，是因为乳酸堆积在肌肉里。”——这个说法流传极广，连不少健身教练都这么说。但反例很清楚：测量表明，剧烈运动后血液和肌肉里的乳酸，在停止运动后\textbf{一小时内}就基本被清除（运走氧化或重新变回糖）。第二天的酸痛姗姗来迟，高峰在运动后一到两天——那时乳酸早就走了。目前更被接受的解释是：不习惯的运动造成肌纤维的微小损伤和随后的修复性炎症，才是延迟性酸痛的来源。乳酸不是废物，更不是“毒素”：它可以被心脏和其他组织当作燃料直接烧掉，运动饮料广告里“清除乳酸”的口号，卖的是安慰剂。当然，酸痛剧烈到反常、或伴随尿液变色等情况，请去看医生，而不是看科普书。
\end{wuqu}

\section{没有活细胞，糖也能“发酵”吗}

\begin{zhengju}
十九世纪中叶，酿酒业已经是欧洲的支柱产业，可没有人知道发酵到底是什么。法国化学家巴斯德用显微镜盯着发酵液看，发现里面总有活的酵母在繁殖；他还发现，把发酵液加热杀死酵母，发酵就停了。于是他宣布：发酵是活细胞的生命活动，“没有生命就没有发酵”。

德国化学家李比希针锋相对。他认为发酵不过是化学物质自行分解，酵母只是碰巧在场。两边吵了几十年，谁也说服不了谁——因为这需要一个决定性的实验：\textbf{把酵母杀死、打碎，看看剩下的“汁”还能不能让糖发酵}。如果能，巴斯德的生命说就站不住；如果不能，李比希就赢了。

难就难在技术上：怎么把细胞彻底弄碎，又不破坏里面的化学物质？1897 年，德国化学家爱德华·毕希纳本来在做一件不相干的事——他想用酵母提取液研究药物的防腐。他把酵母和细沙一起研磨，榨出汁液，为了防腐往里加了大量蔗糖。结果汁液开始冒泡：糖在无细胞的液体里发酵了，酒精和二氧化碳一样不少。毕希纳意识到他撞上了一件大事，重复验证后发表：\textbf{发酵不需要完整的活细胞，活细胞里的化学物质就够了}。

这场争论的结局颇有戏剧性：李比希的“化学说”对了方向，但他以为的那种简单分解并不存在；巴斯德的“生命说”错了，可他坚持的“发酵由精密的生物机制驱动”也没全错——真正干活的是细胞里一整套酶，后来被称为“酵素”（酶这个词的希腊文原意就是“在酵母中”）。接下来的四十多年里，埃姆登、迈耶霍夫、帕尔纳斯等好几代生物化学家像接力一样，用抑制剂卡住某一步、再分离出堆积的中间产物，硬是把十步反应一个个拼凑了出来——这条通路因此也叫“埃姆登–迈耶霍夫–帕尔纳斯途径”。毕希纳拿到了 1907 年的诺贝尔化学奖，而真正的遗产是一个观念：\textbf{生命活动可以在试管里重现、测量、拆解}。生物学和化学之间的那堵墙，从这一杯冒泡的酵母汁开始倒塌。
\end{zhengju}

\begin{shiyan}
\textbf{观察酵母的“呼吸”}（安全，可在家完成）

材料：干酵母一小包（超市有售）、白糖、温水、两个透明塑料瓶、两个气球、记号笔。

步骤：①两个瓶子里各装半瓶温水（手摸不烫，约 $35\,^{\circ}\mathrm{C}$），各加一勺糖，摇匀。②只在其中一个瓶子里加入半包酵母，摇一摇；另一个不加，做对照。③把气球分别套在两个瓶口上，扎紧。④放在温暖处，每隔十分钟观察一次，持续一小时。

观察点：加了酵母的瓶子里，液面是否浮起泡沫？气球是否慢慢鼓起来？对照瓶呢？一小时后取下气球闻闻瓶口（轻轻扇闻），有什么气味？

安全提示：只用食用材料，不使用玻璃密闭容器——气体持续产生，玻璃瓶密封有爆裂风险；实验结束后把液体倒入下水道，气球和瓶子洗净。不要饮用发酵液：它可能混入杂菌，而且实验的乐趣在观察，不在品尝。
\end{shiyan}

\section{这条通路何时够用，何时不够}

任何模型都有边界，糖酵解也不例外。它的强项是\textbf{快}和\textbf{不挑地方}：十步酶全在细胞质里，随叫随到，不需要氧气，不需要细胞器。短跑冲刺的头几十秒，你的肌肉主要靠它撑着；成熟的红细胞干脆连细胞核和线粒体都扔掉了，一辈子只靠糖酵解过活——它要的只是足够运送氧气的 ATP，够用就行。

但它的短板同样明显：每分子葡萄糖只净赚 2 个 ATP，原料利用率低得惊人。长期靠它供能，就像靠烧家具取暖——暖和是暖和，房子很快就没了。这正是它必须和两条“下游通路”配合的原因：有氧气时，丙酮酸和 NADH 进线粒体继续拆（第 53 章）；没氧气时，靠发酵倒电子垃圾维持运转。

医学恰好利用了糖酵解的“脾气”。许多癌细胞有个古怪的嗜好：即使氧气充足，也疯狂吞噬葡萄糖走糖酵解（这个现象叫“瓦氏效应”，原因至今还在研究）。医生让病人注射一种带微弱放射性的葡萄糖类似物，癌细胞吃得越多、信号越强——PET 扫描就这样把肿瘤在片子上“点亮”。本章开篇的那口糖，转身成了侦察兵。顺便提醒一句：这\textbf{不}等于“吃糖喂癌、戒糖治癌”，人体血糖被严格调控，极端戒糖没有证据支持，具体饮食问题请咨询医生。

\section{回到那团面}

现在可以对答案了。酵母没有烧糖——它走的是十步接力：先把葡萄糖拆成丙酮酸，赚下 2 个 ATP；面团深处氧气稀缺，于是启动酒精发酵把 \chem{NAD^{+}} 找回来，顺手扔掉两样电子垃圾：二氧化碳和酒精。二氧化碳被面筋网络兜住，吹出一个个气室，这就是面包的孔；酒精在 $200\,^{\circ}\mathrm{C}$ 左右的烤箱里几乎挥发干净，只留下一丝烘烤香。你摸到的那点温度，正是十步台阶上没来得及被 ATP 接住、以热散掉的那部分能量。

面团没有烫手、没有火苗，因为能量从来不是“一下子放出来的”，而是被酶一级一级拆开、一份一份接走的。火焰和面包的区别，不在放多少能量，而在\textbf{能量释放的速率和去向}——这句话从第 27 章说到现在，在细胞里照样成立。

\section{拆掉只是开始}

回头看这一章，我们其实只干了半件事：把葡萄糖拆到丙酮酸为止。两个丙酮酸和两分子 NADH 还攥着葡萄糖百分之九十以上的能量，细胞绝不会放过它们。在氧气充足的地方，丙酮酸会被请进一座叫线粒体的“发电厂”，被拆到只剩二氧化碳；而 NADH 交出的电子，将沿着一条蛋白质“传送带”逐级下放，最终驱动一台旋转的分子马达成批制造 ATP。电子的终点站着一位我们最熟悉的老朋友——氧气。它为什么非要在终点等着？没有它，传送带会怎样瘫痪？下一章，我们去线粒体里看个究竟。

\begin{tiaozhan}
既然 \chem{NAD^{+}} 这么重要，它凭什么能“接电子”？\chem{NAD^{+}} 的全名叫烟酰胺腺嘌呤二核苷酸，分子里有一个烟酰胺环，恰好能接受一对电子和一个质子，变成 NADH——这本质上是一个氧化还原半反应（第 28 章的语言）。细胞里 \chem{NAD^{+}} 和 NADH 的总和基本固定，两者的比例像一个“电子库存表”：NADH 占上风，说明燃料充足、电子积压，细胞会放慢拆解、转向合成；\chem{NAD^{+}} 占上风，说明库存空虚，糖酵解加速。糖酵解自己的第三步酶（磷酸果糖激酶）就听这个“库存表”指挥：ATP 充足时它直接被 ATP 抑制——产品和原料同源，多产即多停。第 55 章会看到，整个代谢网络就是用这种“下游产品回头管住上游阀门”的办法，把几百条通路拧成一台会自我调节的机器。这些细节跳过不影响主线，但值得知道：调节不是细胞“会做决定”，而是化学亲和力写好的自动反馈。
\end{tiaozhan}

\section*{练一练}

\begin{sikaoti}
\item 画一幅物质流图：从一分子葡萄糖出发，画出“投资 2 ATP、裂成两半、收益 4 ATP 和 2 NADH、得到 2 丙酮酸”的全过程，用箭头标明 ATP 的进和出。算一算：如果把“投资”也算进去，糖酵解的“投资回报率”是多少？
\item 一罐可乐约含 \ud{35}{g} 糖（约 0.2 摩尔）。假如这些葡萄糖全部只走糖酵解，能净得多少摩尔 ATP？对照正文中“完全氧化放能 2870 kJ/mol”，说说为什么只靠糖酵解“烧家具取暖”不是长久之计。
\item 酸奶缸密封不严、进了空气，乳酸菌的发酵效率反而下降——这叫“巴斯德效应”的一个侧面。用本章“\chem{NAD^{+}} 再生”的逻辑解释：为什么氧气一到场，细胞就不那么依赖发酵了？
\item 毕希纳的酵母汁实验为什么能终结“生命说”和“化学说”的争论？如果实验结果是“酵母汁完全不能发酵”，结论又会倒向哪一边？这说明一个判决性实验需要满足什么条件？
\item 有广告说：“本饮料添加左旋乳酸，快速中和运动产生的乳酸，消除疲劳。”请找出这段话里至少两处科学错误，并用本章内容反驳。
\item 成熟的红细胞没有线粒体，只能靠糖酵解供能。请从“红细胞的工作是运送氧气”这个角度推测：红细胞只用糖酵解（不消耗氧气）供能，对它的本职工作有什么好处？
\end{sikaoti}
